%% ==========================================================================
%% DRAFT -- August 12, 2026. Author byline: Daniel Kirtchakov (05oz,
%% halfounce.io; no institutional affiliation; ORCID 0009-0009-5213-4098).
%%
%% Part K (wedge v2). Supersedes the core result of Part H
%% (doi:10.5281/zenodo.21895825): the two-sided bracket L <= P_L <= U
%% collapses to the exact rational P_L, and the weight-7 count wall named in
%% Part H, section 6 falls at d=5. Framing constraints honoured (PROTOCOL
%% section 11): the note opens with the question Part H itself asked (its
%% Question 7.1), states the identity as real mathematics with proof, reports
%% the honest trade (the C(m,w) wall is exchanged for a 2^n wall; d=7 remains
%% out of reach on a laptop and the note says so precisely), cites the
%% computation once in a Certification paragraph, and closes with the
%% questions the computation opens. The novelty claim is scoped to the
%% intersection recorded in SWEEP-RECORD-WEDGE2-2026-08-12.md; the character
%% sum is acknowledged as MacWilliams-type and classical in lineage.
%% ==========================================================================
\documentclass[11pt]{amsart}

\usepackage[margin=1in]{geometry}
\usepackage{amsmath,amssymb}
\usepackage{booktabs}
\usepackage{array}
\usepackage[colorlinks=true,allcolors=blue]{hyperref}

\emergencystretch=3em
\tolerance=1200
\hbadness=4000

\newcommand{\PL}{P_{L}}
\newcommand{\sig}{\sigma}
\newcommand{\obs}{\mathrm{obs}}
\newcommand{\Prb}{\mathbb{P}}
\newcommand{\Zt}{\mathbb{F}_2}
\newcommand{\FWHT}{\widehat{\phantom{x}}}

\theoremstyle{plain}
\newtheorem{theorem}{Theorem}[section]
\newtheorem{lemma}[theorem]{Lemma}
\newtheorem{proposition}[theorem]{Proposition}
\newtheorem{corollary}[theorem]{Corollary}
\theoremstyle{definition}
\newtheorem{definition}[theorem]{Definition}
\theoremstyle{remark}
\newtheorem{remark}[theorem]{Remark}
\newtheorem{question}[theorem]{Question}

\begin{document}

\title[The exact logical error probability]{The exact logical error
probability of the rotated surface code under a lookup-table decoder:
syndrome-space character sums collapse a certified bracket to a single
rational, replayable in the standard library}

\author{Daniel Kirtchakov}
\address{Independent researcher (\texttt{05oz}), Half Ounce Research; no
institutional affiliation. ORCID
\href{https://orcid.org/0009-0009-5213-4098}{0009-0009-5213-4098}.}
\email{daniel@halfounce.io}
\urladdr{https://halfounce.io}

\date{August 12, 2026}

\thanks{\textbf{Computation and authorship.} All derivation, implementation,
verification, and drafting in this work were produced by \textbf{Claude}
(Anthropic), directed by the author, on a single Apple laptop. The public
artifact is a certificate together with a checker that imports
only the Python standard library and re-derives every certified quantity
from the mechanism list alone. The checker is not code-independent of the
private engine that produced the certificate. Of the 238 executable lines of
\texttt{check\_wedge2.py}, 52 appear verbatim in that engine, among them an
eleven-line byte-identical run, lines 116 to 126, which is the untruncated
breadth-first search that rebuilds the coset-leader decoder; the
detector-mask histograms that feed the transforms, the binning of mechanisms
by probability, the scaled character-power tables, and the construction of
the counting polynomials $g_w$ are shared as well. A passing check may
therefore not be read, on its own, as an independent re-derivation of the
decoder or of the $A_w$ and $N_w$ spectra.\endgraf
What is original to the checker is the transform routine itself, which has
only two lines in common with the engine, \texttt{h = 1} and \texttt{return
a}; the exact integer inversion and its integrality guards; the invariants
$\sum_w A_w = 2^{m-1}$ and $\sum_w N_w = 2^{m-n-1}$; the mechanism-list hash
recomputation; every comparison against the certificate; and the
fail-and-exit reporting. The tamper controls rest on that layer and are
unaffected.\endgraf
At $d=3$ the exact $\PL$ has a check that does not depend on this checker at
all. Both $d=3$ certificates carry the same mechanism list as Part H's ---
the SHA-256 of the canonical list agrees --- and an implementation written
against the definitions of Part H, sharing no routine with the Part H
checker or its engine, re-derives that bracket end to end: the coset-leader
decoder at all 256 syndromes, the uncorrectable counts 55, 690 and 4,078 at
weights 2, 3 and 4, all 4,823 uncorrectable sets, and $L$, $T$, $U$ exactly.
The exact $\PL$ certified here lies strictly inside that independently
re-derived bracket at both $d=3$ operating points, and those three counts
agree with the low-weight end of the $A_w$ spectrum certified here. At $d=5$
there is no such check: the independent implementation reproduces Part H's
tail $T$ there but not its lower bound $L$, so containment of $\PL$ in the
Part H bracket at $d=5$ confirms nothing independently, and no part of the
$d=5$ $A_w$ or $N_w$ spectrum has been independently re-derived. The
detector error models are those of Part H
\cite{PartH}, generated there with Stim 1.16.0 (Gidney), which is a generator
and is \emph{not} in the trusted base (\S5).}

\thanks{\textbf{Prior-art record.} Novelty was swept on August 12, 2026
against arXiv and the web; the dated record ships with the artifacts as
\texttt{SWEEP-RECORD-WEDGE2-2026-08-12.md}. The character-sum technique is
MacWilliams-type and classical in lineage; the claim is scoped to the
intersection stated in \S3.}

\begin{abstract}
Part H of this program \cite{PartH} published machine-checkable two-sided
rational brackets $L \le \PL \le U$ on the logical error probability of the
distance-3 and distance-5 rotated surface codes, one round, circuit-level
depolarizing noise, under a fixed lookup-table (coset-leader) decoder --- and
named its own frontier: exact re-verification of the uncorrectable-set counts
at weight 7 exceeds a pure-Python checker, because the shipped method
enumerates $\binom{77}{7} = 2{,}404{,}808{,}340$ fault sets. This note
removes both the truncation and the wall at $d \le 5$. A syndrome-space
character sum (a MacWilliams-type Walsh--Hadamard argument) yields
\emph{every} uncorrectable count $A_w$, $w = 0, \dots, m$, in a single
$O(2^n n)$ pass over the $2^n$ syndromes --- no per-weight enumeration ---
and its probability-weighted variant yields the \emph{exact} rational $\PL$,
eliminating the bracket altogether. At $d=5$ this delivers the previously
unreachable $A_7 = 832{,}441{,}445$, the full spectrum through $A_{77} = 1$,
the undetectable-logical spectrum $N_w$ (re-deriving the circuit-level
distance~5), and exact dyadic rationals $\PL$ at both published operating
points, each lying strictly inside the corresponding Part H bracket --- at
$0.36$ and $0.37$ of the bracket widths, never near an endpoint. A
standard-library Python checker re-derives all of it from the mechanism list
alone in about two minutes at $d=5$ within $0.85$ GB, and rejects eight
classes of tampered certificate. The honest trade: the
$\binom{m}{w}$ enumeration wall is exchanged for a $2^n$ syndrome-space wall,
so one-round distance-7 codes ($n \gtrsim 48$ detectors) remain out of reach
on a laptop by this method, and we state precisely why. The exact values
supersede the Part H brackets; the brackets survive as the independent check
the exact values passed.
\end{abstract}

\maketitle

\section{The question}\label{sec:question}

Part H \cite{PartH} asked whether the logical error probability $\PL$ of a
quantum error-correcting code, in the deep sub-threshold regime that defeats
direct Monte Carlo sampling \cite{MWB}, can be pinned by a certificate a
skeptic replays with nothing but a standard-library Python interpreter. Its
answer was a two-sided exact-rational bracket $L \le \PL \le U$, built from
exhaustively enumerated uncorrectable-set counts $A_w$ up to a truncation
weight $W_{\max}$, with the width $U - L$ a Poisson-binomial tail. Its
\S6 then located the method's own frontier, and its Question 7.1 asked:

\begin{quote}
\emph{Is there a proof-carrying counter for the exact weighted enumeration at
$W_{\max} = 7$ --- a CPOG-style weighted model-counting certificate, or an
exact locality-exploiting convolution / transfer-matrix --- that a
standard-library checker can re-verify in minutes rather than hours?}
\end{quote}

This note answers that question affirmatively at $d \le 5$, and with a
stronger object than the question requested. The answer is not a faster
enumeration and not a tighter bracket: a character sum over the syndrome
space computes all counts $A_w$ simultaneously, with cost independent of the
weight, and its probability-weighted form computes $\PL$ \emph{exactly} ---
as a single reduced rational --- so that the truncation weight, the tail, and
the bracket disappear from the certified statement. The Part H brackets are
retained by the new certificates in one role only: as published, independent
enclosures that the exact values are verified to satisfy. Everything else
about the Part H trust architecture --- the mechanism list as trust root, the
fixed decoder, the standard-library checker, the private generator --- is
unchanged, and the new checker re-verifies in two minutes, at all weights at
once, what the enumeration checker could not reach at weight 7 in hours.

\section{The identity}\label{sec:identity}

\textbf{Model.} Fix a detector error model (DEM): $n$ detectors, one logical
observable, $m$ independent fault mechanisms; mechanism $j$ carries a
detector mask $d_j \in \Zt^n$, an observable bit $o_j \in \Zt$, and a
probability $p_j \in (0,1)$, an exact dyadic rational (the IEEE-754 double
emitted by Stim). Write $v_j = (d_j, o_j) \in \Zt^{n+1}$. A fault
configuration $F \subseteq [m]$ has syndrome $\sig(F) = \bigoplus_{j \in F}
d_j$, observable flip $\obs(F) = \bigoplus_{j \in F} o_j$, weight $|F|$, and
probability $\Prb(F) = \prod_{j \in F} p_j \prod_{j \notin F} (1 - p_j)$.

\textbf{Decoder.} The decoder $D : \Zt^n \to \Zt$ is the coset-leader lookup
table of Part H: a full breadth-first search over the syndrome space from the
zero syndrome, expanding by the mechanism masks in canonical index order,
first assignment winning; $D(s)$ is the observable flip of the first-found
minimum-cardinality configuration reaching $s$. In the two DEMs certified
here the masks span $\Zt^n$, so every syndrome is assigned (the checker
verifies this), and $F$ is uncorrectable iff $D(\sig(F)) \ne \obs(F)$. On
syndromes of BFS depth $\le W_{\max}$ this table coincides with the table
Part H certifies --- built there in full at $d=3$ and truncated at depth
$W_{\max}$ at $d=5$ --- so the decoder is the same decoder.

\textbf{The counts, all at once.} For $u \in \Zt^{n+1}$ let $I(u)$ be the
failure indicator: $I(s, b) = 1$ iff $b \ne D(s)$. Let $A_w$ be the number of
uncorrectable $F$ with $|F| = w$, and let
$N_w = \#\{F : |F| = w,\ \sig(F) = 0,\ \obs(F) = 1\}$ count the undetectable
logical operators, so that the circuit-level distance is $\min\{w : N_w >
0\}$. For $y \in \Zt^{n+1}$ write $\langle y, u\rangle$ for the standard
pairing, $\widehat{f}(y) = \sum_u (-1)^{\langle y,u \rangle} f(u)$ for the
(unnormalised) Walsh--Hadamard transform, and
\[
a(y) \;=\; \#\{\, j : \langle y, v_j\rangle = 0 \,\}, \qquad
g_w(k) \;=\; [t^w]\,(1+t)^k (1-t)^{m-k} .
\]

\begin{theorem}[Counting identity]\label{thm:star}
$\displaystyle A_w \;=\; \frac{1}{2^{\,n+1}} \sum_{y \in \Zt^{n+1}}
\widehat{I}(y)\, g_w(a(y))$, and the same identity with $\widehat{I}$
replaced by the transform of the indicator of $(0,1)$-cosets gives
$\displaystyle N_w = \frac{1}{2^{\,n+1}} \sum_{y \in \Zt^{n+1}} (-1)^{y_b}\,
g_w(a(y))$, where $y = (y_s, y_b)$, $y_s \in \Zt^n$, $y_b \in \Zt$.
\end{theorem}

\begin{proof}
Let $N_w(u) = \#\{F : |F| = w,\ \bigoplus_{j \in F} v_j = u\}$. In the group
algebra of $\Zt^{n+1}$, $\sum_{w,u} N_w(u)\, t^w e_u = \prod_j (1 + t\,
e_{v_j})$, and the character $\chi_y(e_u) = (-1)^{\langle y,u\rangle}$ maps
the product to $\prod_j (1 + t(-1)^{\langle y, v_j\rangle}) =
(1+t)^{a(y)} (1-t)^{m - a(y)}$. Fourier inversion gives $N_w(u) =
2^{-(n+1)} \sum_y (-1)^{\langle y,u\rangle} g_w(a(y))$, and $A_w = \sum_u
I(u) N_w(u) = 2^{-(n+1)} \sum_y \widehat{I}(y)\, g_w(a(y))$ by exchanging
the finite sums. Taking $u = (0,1)$ instead of summing against $I$ gives the
$N_w$ statement, since $(-1)^{\langle y, (0,1)\rangle} = (-1)^{y_b}$.
\end{proof}

\begin{theorem}[Exact probability identity]\label{thm:dia}
$\displaystyle \PL \;=\; \frac{1}{2^{\,n+1}} \sum_{y \in \Zt^{n+1}}
\widehat{I}(y) \prod_{j\,:\, \langle y, v_j\rangle = 1} (1 - 2p_j)$.
\end{theorem}

\begin{proof}
Identical character argument applied to $\prod_j \bigl((1-p_j)\,e_0 + p_j\,
e_{v_j}\bigr)$: the coefficient of $e_u$ is the total probability of
$\{F : \bigoplus_{j\in F} v_j = u\}$, and $\chi_y$ maps each factor to
$(1-p_j) + p_j(-1)^{\langle y, v_j\rangle}$, which is $1$ when
$\langle y, v_j \rangle = 0$ and $1 - 2p_j$ otherwise. Summing the inverted
coefficients against $I$ gives $\PL = \sum_u I(u) \cdot [\text{coefficient}]$
as claimed.
\end{proof}

Splitting $y = (y_s, y_b)$ makes both identities computable by transforms of
$n$-bit arrays: with $c(s) = 1$ (all syndromes reachable), $e(s) =
(-1)^{D(s)}$, one finds $\widehat{I}(y_s, 0) = 2^n[y_s = 0]$ and
$\widehat{I}(y_s, 1) = -\widehat{e}(y_s)$; and $a(y)$ becomes $a_0(y_s) =
\#\{j : \langle y_s, d_j\rangle = 0\}$ resp.\ $a_1(y_s) = \#\{j : \langle
y_s, d_j\rangle = o_j\}$, each obtainable for all $y_s$ from one transform of
a sparse mask histogram. Four transforms of length $2^n$ therefore produce
every $A_w$ and every $N_w$; the total cost is $O(2^n n)$ arithmetic
operations on integers of magnitude at most $2^n$, plus the $O(2^n m)$ BFS
that builds $D$ --- \emph{independent of $w$}, which is what removes the
$\binom{m}{w}$ wall. For Theorem~\ref{thm:dia} the product
$\prod_{\langle y, v_j\rangle = 1} (1 - 2p_j)$ depends only on how many
mechanisms of each distinct probability class have $\langle y, v_j \rangle =
1$ (17 classes at $d=5$, 14 at $d=3$), so the exact big-integer assembly
groups the $2^n$ characters by a class-count key; all denominators are
powers of two, so the common-denominator scaling is a bit shift and the
result is a dyadic rational, reduced at the end. Cancellation is severe ---
the sum is $2^n$ minus a nearly-equal character sum --- which is exactly why
the arithmetic is exact integers throughout; a floating-point evaluation of
Theorem~\ref{thm:dia} would be meaningless at these magnitudes.

Two closed-form invariants gate the whole computation.

\begin{proposition}\label{prop:invariants}
If some $F_0$ satisfies $\sig(F_0) = 0$, $\obs(F_0) = 1$ (true here: $N_5 =
111 > 0$ at $d = 5$, $N_3 = 9 > 0$ at $d=3$), then $\sum_{w} A_w = 2^{m-1}$.
If moreover the $(d_j, o_j)$ span $\Zt^{n+1}$ (true here), then $\sum_w N_w
= 2^{\,m-n-1}$.
\end{proposition}

\begin{proof}
$F \mapsto F \,\triangle\, F_0$ is an involution on subsets fixing
$\sig$ and flipping $\obs$; since $D(\sig)$ is fixed, exactly one of each
pair is uncorrectable, giving $2^m/2$. The second claim: $F \mapsto
(\sig(F), \obs(F))$ is a surjective homomorphism $(\Zt^m, \triangle) \to
\Zt^{n+1}$, so each fiber has $2^{m-(n+1)}$ elements.
\end{proof}

\section{Where this sits}\label{sec:placement}

The character sum is not new mathematics: computing the probability of
correct syndrome decoding from coset weight enumerators via MacWilliams-type
transforms is classical coding theory \cite{JP}. What the classical setting
lacks is everything the certificate needs: there, the channel is symmetric
with one parameter and the error positions are the code coordinates; here,
$m > n$ heterogeneous dyadic mechanisms carry an extra observable bit, the
decoder is a specific tabulated function fixed by a BFS convention, and the
output is an exact rational that an independent standard-library checker
re-derives, hash-bound to a mechanism list. The neighbouring modern
literature splits as in Part H: estimators return values with error bars
\cite{MWB, BV}; tensor-network maximum-likelihood computations
\cite{BSV, DP} are numerically exact floating-point contractions with no
certified output; closed-form and enumerator-based analyses \cite{Regev,
FVC} are approximations or code-level bounds, not decoder-specific exact
values; recent uses of Walsh--Hadamard relations on DEMs
\cite{Willow, CAHR} estimate or reconstruct the model from data rather than
certify anything about a decoder; formal verifiers \cite{LeanQEC, VeriQEC}
machine-check Boolean error-correction conditions, not probabilities. The
dated sweep (2026-08-12) found no prior exact certified logical error
probability of a stated decoder on a circuit-level DEM, by transform methods
or otherwise, with independent standard-library re-verification; we claim
that intersection and nothing broader, and this note supersedes our own
Part H bracket \cite{PartH} rather than any external result.

\section{Results}\label{sec:results}

The two DEMs are byte-identical to Part H's (same mechanism lists, same
SHA-256 hashes): $d = 3$ with $n = 8$, $m = 23$; $d = 5$ with $n = 24$,
$m = 77$. All counts below are $p$-independent; the exact $\PL$ is computed
at the two published operating points of each code.

\textbf{The spectra.} At $d = 5$ the transform yields in one pass the full
uncorrectable spectrum $A_0, \dots, A_{77}$, whose head is
\[
A_3 = 2728, \quad A_4 = 154{,}394, \quad A_5 = 4{,}057{,}999, \quad
A_6 = 67{,}711{,}204, \quad \mathbf{A_7 = 832{,}441{,}445},
\]
continuing through $A_{77} = 1$ (the full-support configuration). $A_3,
A_4, A_5$ equal the Part H certified counts; $A_6$ equals the value of the
optional $W_{\max} = 6$ certificate of Part H; $A_7$ and beyond are new ---
weight 7 is precisely where Part H's enumeration hit its wall
($\binom{77}{7} \approx 2.4 \times 10^9$ subsets, hours of pure-Python
re-verification). The undetectable-logical spectrum begins $N_5 = 111$,
$N_6 = 786$, $N_7 = 3706$, re-deriving circuit-level distance $5$ by
transform rather than enumeration. Both closed-form invariants hold to the
digit: $\sum_w A_w = 2^{76}$ and $\sum_w N_w = 2^{52}$ (at $d=3$: $2^{22}$
and $2^{14}$). At $d = 3$ the head $A_2 = 55$, $A_3 = 690$, $A_4 = 4078$
matches Part H, and $A_5 = 16{,}467$ was confirmed by direct enumeration.

\textbf{The exact values.} Table~\ref{tab:exact} is the core supersession:
the exact $\PL$ at all four published operating points, with the Part H
bracket it collapses. Every exact value lies \emph{strictly inside} its
bracket. What the shipped checker gates is the containment $L \le \PL \le U$,
verified in exact rational arithmetic on the certificate's own numerators and
denominators; strictness is not part of that gate, but follows from the same
three numbers, which place every exact value between $0.35$ and $0.52$ of the
bracket width (last column of Table~\ref{tab:exact}). Had the containment
failed the release would have stopped, since it would have falsified Part H.

\begin{table}[ht]
\centering
\caption{Exact $\PL$ (decimal shadows of exact dyadic rationals; the
certificates carry full numerators and denominators, e.g.\ 1026/1030 digits
at $d{=}5$, $p{=}10^{-3}$) against the Part H brackets $[L, U]$. ``Position''
is $(\PL - L)/(U - L)$.}
\label{tab:exact}
\begin{tabular}{lllll}
\toprule
code, $p$ & exact $\PL$ & Part H $L$ & Part H $U$ & position \\
\midrule
$d{=}3$, $10^{-3}$ & 3.358965645735e-4 & 3.3589593231e-4 & 3.3589715988e-4 & 0.515 \\
$d{=}3$, $10^{-2}$ & 2.714610979781e-2 & 2.7101553370e-2 & 2.7188139126e-2 & 0.515 \\
$d{=}5$, $10^{-3}$ & 2.613858119750e-5 & 2.6135024167e-5 & 2.6144956889e-5 & 0.358 \\
$d{=}5$, $10^{-2}$ & 1.734152068512e-2 & 1.6118428758e-2 & 1.9426205101e-2 & 0.370 \\
\bottomrule
\end{tabular}
\end{table}

Three further consistency checks. The $d{=}5$, $p = 10^{-3}$ exact value
lies inside the tighter optional $W_{\max}{=}6$ bracket of Part H,
$[2.6138502142\mathrm{e}{-5},\, 2.6138690261\mathrm{e}{-5}]$, at position
$0.42$ --- a check about $53\times$ more stringent than the released
$W_{\max}{=}5$ bracket. The exact values at both $d=5$ points lie inside the
$10^7$-shot Monte Carlo 95\% confidence intervals recorded in Part H. And at
$p = 10^{-2}$, $d = 5$ --- the regime where Part H's bracket was honestly
reported \textbf{DEAD}, $20.5\times$ wider than Monte Carlo --- the exact
value $1.734152\mathrm{e}{-2}$ settles what the bracket could not resolve:
the bracket is now beaten not by more shots but by eliminating the
truncation altogether. The supersession is uniform: where the bracket was
LIVE it is replaced by the exact value; where it was DEAD it is replaced by
the exact value; nothing bracketed remains.

\textbf{Independent confirmations (pre-registered standard).} Because the
headline numbers pass beyond what Part H could re-verify, each was confirmed
by at least one computational route sharing no arithmetic with the transform
engine. At $d = 3$, \emph{three} independent exact routes agree digit-for-digit
on $\PL$ at both operating points: the class-binned character sum
(Theorem~\ref{thm:dia}), a syndrome-space signed convolution reconstructed
over 25-bit primes by the Chinese remainder theorem, and a Gray-code full
enumeration of all $2^{23}$ fault configurations in exact integer
arithmetic. At $d = 5$, where full enumeration is impossible ($2^{77}$
configurations), the character-sum values were checked against the
independent convolution route modulo 24 distinct 25-bit primes ---
$\approx 600$ bits of agreement --- at both operating points, and the count
spectrum was checked against every Part H certified count, the $W_{\max}=6$
engine value, both invariants of Proposition~\ref{prop:invariants}, and Part
H's Monte Carlo. Before this build, the scoping run had also cross-checked
$A_7$ by a truncated-BFS Monte Carlo estimator ($z = 0.19$).

\textbf{Certification.} The permanent record is four certificate JSONs ---
the mechanism list and hash (identical to Part H), the full $A_w$ and $N_w$
spectra, the circuit-level distance, both invariants, the exact $\PL$
numerator and denominator, and the cited Part H bracket --- together with
one standard-library checker, \texttt{check\_wedge2.py}, that re-derives
every certified quantity from the mechanism list alone: it rebuilds the full
BFS decoder, verifies spanning, runs the Walsh--Hadamard transforms on
\texttt{array('q')} buffers, assembles the spectra and the exact rational
$\PL$ in big-integer arithmetic, and verifies $L \le \PL \le U$ against the
embedded Part H bracket. It imports only \texttt{hashlib}, \texttt{json},
\texttt{sys}, \texttt{array}, \texttt{fractions}; no signals, subprocesses,
network, or wall-clock. Runtime: $0.15$ s at $d=3$; at $d=5$, $129$ s and
$0.79$--$0.84$ GB peak resident memory per certificate on the build laptop
(measured with \texttt{/usr/bin/time -l}) --- for comparison, Part H's
enumeration checker took a comparable $34$ s to certify weights ${\le}5$
only, ten-plus minutes for ${\le}6$, and hours for ${\le}7$; this checker
certifies all $78$ weights \emph{and} the exact $\PL$ in two minutes. A
tamper battery of eight corruptions --- a flipped count, a perturbed
probability with and without a recomputed hash, an altered $\PL$ digit, a
swapped observable bit, an overstated distance, a corrupted $N_w$, a shifted
bracket --- is rejected gate-by-gate on the $d{=}3$ certificate, where the
battery replays in under a second (\texttt{tamper\_demo\_w2.py}; the
corruptions are certificate-generic and the script runs against any of the
four). The
identity itself ships with a standard-library self-test
(\texttt{identity\_selftest.py}) that verifies both theorems against brute
force on random small DEMs, including non-spanning mask sets, in exact
\texttt{Fraction} arithmetic. The transform engine, like Part H's
enumeration engine, is private and is not needed to check anything.

\section{What is certified, and what is trusted}\label{sec:trust}

The trust boundary is unchanged from Part H \cite{PartH}, and is restated
plainly. \textbf{Machine-checked:} everything in \S4 --- spectra,
invariants, distance, the exact $\PL$, bracket containment --- is re-derived
by the standard-library checker from the certificate's mechanism list, which
is hash-pinned and byte-identical to Part H's. \textbf{Trusted:} (1) the
binding of that mechanism list to Stim's DEM at the stated $p$ lives in the
private generator and cannot be re-checked from the public artifact; (2) the
independent-mechanism DEM is Stim's standard approximation of the physical
depolarizing circuit, and the certified $\PL$ is the DEM's, not the
circuit's --- Part H measured the gap at $d=5$, $p=10^{-3}$ (a raw-circuit
Monte Carlo point estimate about $1.17$ standard errors above the DEM
bracket) and that caveat carries over verbatim; (3) CPython, the OS, the
hardware; (4) SHA-256, for bookkeeping only. One trusted item is
\emph{removed}: the depth-truncation soundness argument of Part H (that a
depth-$W_{\max}$ BFS agrees with the full table on low-weight syndromes) is
no longer needed, because the checker builds the full table.

\section{The wall that replaces the wall}\label{sec:wall}

Part H's frontier was the enumeration $\binom{m}{w}$; this note's method
replaces it with the syndrome space $2^n$. The exchange is decisively
favourable at $d = 5$ --- $2^{24} \approx 1.7 \times 10^7$ versus
$\binom{77}{7} \approx 2.4 \times 10^9$, with all weights for free --- and
unfavourable beyond. Every step of the method --- the full BFS decoder, the
$2^n$-length transforms, the class-binned assembly --- materialises arrays
of length $2^n$. One round of the distance-7 rotated surface code has $n
\approx 48$ detectors: $2^{48} \approx 2.8 \times 10^{14}$ syndromes, seven
orders of magnitude beyond what fits in laptop memory, before any
arithmetic. The obstruction to doing better is specific and worth naming:
the masks are local (at $d=5$, every mechanism flips at most two detectors),
so the mask-side transforms would factor beautifully over any tensor
decomposition of the syndrome space --- but the decoder $D(s)$ is a
\emph{global} minimum-cardinality object, defined by a BFS over the whole
space, and it is $\widehat{e}$, the transform of $(-1)^{D(s)}$, that the
identities cannot do without. A method that reaches $d = 7$ must either
localise this global object (a frontier transfer-matrix that carries
coset-leader information without materialising $2^n$ states) or abandon the
lookup decoder for one with local structure. We state this as the honest
limit: \emph{by the method of this note, on a laptop, $d = 7$ circuit-level
is out of reach, and no amount of engineering of the present checker
changes that.}

\section{Questions}\label{sec:questions}

\begin{question}
(Part H's Question 7.2, sharpened.) Is there an exact representation of the
coset-leader decoder --- or of any explicitly specified decoder of
comparable quality --- whose failure transform $\widehat{e}$ admits
evaluation without materialising $2^n$ states, opening $d = 7$ at one round
to a laptop-checkable exact certificate?
\end{question}

\begin{question}\label{q:ml}
The signed syndrome measure $q(s) = \sum_{F : \sig(F) = s} \Prb(F)
(-1)^{\obs(F)}$, computable exactly by the convolution route used here as a
cross-check, determines the maximum-likelihood decoder for the observable:
$D_{\mathrm{ML}}(s) = [q(s) < 0]$. The identities of \S2 apply verbatim to
\emph{any} tabulated decoder. Can the exact $\PL$ of the maximum-likelihood
decoder --- the code's true operational figure, Part H's Question 7.4 ---
be certified at $d = 5$ by combining the two routes, with the sign
tabulation itself made checker-re-derivable?
\end{question}

\begin{question}
Multiple rounds multiply the detector count, so the $2^n$ wall arrives
immediately even at $d = 3$. Is there a round-by-round factorisation of the
character sum --- the time direction is where the DEM \emph{is} local ---
that trades $2^n$ for per-round transforms of the single-round syndrome
space?
\end{question}

\begin{question}
Part H's Question 7.3 (certify the DEM-to-circuit gap) is untouched by this
note and remains open; the exact DEM value makes it sharper, since the gap
is now the only distance between the certificate and the physical circuit.
\end{question}

\section*{Acknowledgments}

This note upgrades Part H of this program and answers its Question 7.1; the
question, the decoder convention, the DEMs, the operating points, and the
kill-test record are all Part H's, and the supersession is of our own
published bracket, not of any external result. The framing that
sub-threshold logical error rates defeat direct sampling is Mullan,
Weippert, and Brown's \cite{MWB}. The lookup-table decoder is Tomita and
Svore's \cite{TS}. The coset-enumerator lineage of the character sum is
classical \cite{JP}. The detector error models were produced with Stim
(Craig Gidney), used as an untrusted generator. The computation and drafting
were AI-assisted as stated in the first footnote; the note's shape --- open
with the question, cite the computation once, close with the questions it
opens --- follows the program's standing constraint.

\begin{thebibliography}{99}

\bibitem[PartH]{PartH} D. Kirtchakov, \emph{Certified sub-threshold logical
error rates: exact uncorrectable-set counts and a two-sided rational bracket
for the rotated surface code, replayable in the standard library}, Certify
v0.8.0 (Part H), Zenodo (2026), doi:10.5281/zenodo.21895825.

\bibitem[MWB]{MWB} M. Mullan, M. Weippert and W. Brown, \emph{Improved
methods for determining quantum error correcting code performance and fault
tolerance}, arXiv:2607.27153 (2026).

\bibitem[BV]{BV} S. Bravyi and A. Vargo, \emph{Simulation of rare events in
quantum error correction}, Phys. Rev. A \textbf{88} (2013), 062308.

\bibitem[BSV]{BSV} S. Bravyi, M. Suchara and A. Vargo, \emph{Efficient
algorithms for maximum likelihood decoding in the surface code}, Phys. Rev.
A \textbf{90} (2014), 032326; arXiv:1405.4883.

\bibitem[DP]{DP} A. S. Darmawan and D. Poulin, \emph{Tensor-network
simulations of the surface code under realistic noise}, Phys. Rev. Lett.
\textbf{119} (2017), 040502; arXiv:1607.06460.

\bibitem[JP]{JP} R. Jurrius and R. Pellikaan, \emph{Codes, arrangements and
matroids}, in: Algebraic Geometry Modeling in Information Theory, World
Scientific (2013); and standard references on coset weight enumerators and
the probability of correct syndrome decoding.

\bibitem[Regev]{Regev} G. Regev, G. Dilley, T. Nutaro, R. Delgado and
R. Bennink, \emph{Closed form logical error rate approximations for surface
codes}, arXiv:2605.03054 (2026).

\bibitem[FVC]{FVC} D. Forlivesi, L. Valentini and M. Chiani,
\emph{Performance analysis of quantum error-correcting codes via MacWilliams
identities}, Quantum \textbf{9} (2025), 1950; arXiv:2305.01301.

\bibitem[Willow]{Willow} \emph{Estimating detector error models on Google's
Willow}, arXiv:2512.10814 (2026).

\bibitem[CAHR]{CAHR} \emph{Reconstruction of detector error model for
quantum error correction}, arXiv:2606.16288 (2026).

\bibitem[LeanQEC]{LeanQEC} K. Ehatamm, S. Lee, Y. Wu and R. Tao,
\emph{Lean-QEC: machine-checked quantum error-correcting code distances via
a verified SAT reduction}, arXiv:2605.16523 (2026).

\bibitem[VeriQEC]{VeriQEC} C. Huang et al., \emph{Veri-QEC: verifying error
correction and detection conditions for quantum error correcting codes},
Proc. ACM Program. Lang. (2025), DOI 10.1145/3729293.

\bibitem[TS]{TS} Y. Tomita and K. M. Svore, \emph{Low-distance surface codes
under realistic quantum noise}, Phys. Rev. A \textbf{90} (2014), 062320.

\end{thebibliography}

\end{document}
